% ============================================================
% SECTION 01 — THEORY: ENTANGLEMENT ENTROPY–GRAVITY COUPLING
% Repository: Monette Research — Entropic Gravity & PEIG Framework
% Part of: 01_PHYSICS_CORE/A_ENTROPIC_GRAVITY/master_paper.tex
% Status: Peer-review draft — Phase 2 corrections applied
% Last Updated: March 2026
% ============================================================
%
% PHASE 2 CORRECTIONS APPLIED IN THIS FILE:
%   [C1] kappa-tilde = -1/4 reframed as "first-principles estimate under
%        holographic regulator hypothesis," not "derived result"
%   [C2] Screening factor formula variables explicitly defined
%   [C3] Negentropy/entropy conflation explicitly resolved
%   [C4] Conservation equation tension with precision tests noted
% ============================================================

\section{Theory: Entanglement Entropy--Gravity Coupling}
\label{sec:theory}

\subsection{Physical Motivation}
\label{subsec:motivation}

The central question of this framework is deceptively simple: if information
is physical, and if entropy is the fundamental quantity connecting quantum
mechanics to thermodynamics, then does entropy \emph{source} spacetime
curvature in the same way mass-energy does?

Jacobson's 1995 derivation~\cite{jacobson1995} established that Einstein's
field equations are not a fundamental law of nature but a thermodynamic
equation of state: they follow from applying the Clausius relation
$\delta Q = T\,dS$ to local Rindler horizons, where the entropy is
Bekenstein-Hawking horizon entropy and the temperature is the Unruh
temperature of an accelerated observer. In Jacobson's framework, spacetime
geometry is the thermodynamic equilibrium state of some underlying
microscopic theory.

If this is correct, then any physical source of entropy --- not just
horizon entropy --- should, in principle, back-react on spacetime geometry.
The question becomes quantitative: \emph{how strongly} does laboratory-scale
quantum entanglement entropy couple to spacetime curvature, and is this
coupling measurable with existing technology?

This section derives the modified Einstein equation incorporating
entanglement entropy as a source term, establishes its dimensional
consistency, and derives the ideal value of the coupling constant
$\kappatilde$ from first principles. The critical extrapolation from
horizon physics to laboratory physics is treated honestly in
Section~\ref{sec:extrapolation}: it is a physical hypothesis, not a
mathematical derivation, and it is falsifiable precisely because it makes
quantitative predictions.

% --------------------------------------------------------
\subsection{The Modified Einstein Equation}
\label{subsec:modified_einstein}
% --------------------------------------------------------

The coupling between entanglement entropy density and spacetime geometry
is expressed through the following modification of Einstein's field
equations:

\begin{equation}
    \boxed{
    G_{\mu\nu} = 8\pi G\,T_{\mu\nu}
    + \kappatilde\,\frac{c^4}{k_B\lnTwo}\,\Sent\,g_{\mu\nu}
    }
    \label{eq:modified_einstein}
\end{equation}

\noindent where every quantity is defined as follows:

\begin{itemize}
    \item $G_{\mu\nu} = R_{\mu\nu} - \frac{1}{2}R\,g_{\mu\nu}$ is the
    \textbf{Einstein tensor} --- a measure of spacetime curvature at each
    point. It is constructed from the Ricci tensor $R_{\mu\nu}$ (which
    encodes local curvature) and the Ricci scalar $R = g^{\mu\nu}R_{\mu\nu}$
    (the trace of curvature). Units: m$^{-2}$.

    \item $G = 6.674\times10^{-11}$~m$^3\cdot$kg$^{-1}\cdot$s$^{-2}$ is
    Newton's \textbf{gravitational constant}.

    \item $T_{\mu\nu}$ is the standard \textbf{stress-energy tensor} of
    ordinary matter and radiation --- the source term in the original Einstein
    equations. Its time-time component $T_{00}$ equals energy density; its
    space-space components encode pressure. Units:
    kg$\cdot$m$^{-1}\cdot$s$^{-2}$ (equivalently, Pa = N/m$^2$).

    \item $\kappatilde$ is a \textbf{dimensionless coupling constant} --- the
    central parameter of this framework. Its sign determines whether the
    entanglement entropy contribution produces attractive ($\kappatilde > 0$)
    or repulsive ($\kappatilde < 0$) curvature. Its magnitude determines the
    strength of the coupling. This paper derives the ideal value
    $\kappatilde = -1/4$ and proposes experiments to measure it.

    \item $c = 2.998\times10^8$~m/s is the speed of light. The factor $c^4$
    has units m$^4\cdot$s$^{-4}$, and combined with $G$ in the denominator
    gives $c^4/G$ the units of force (N = kg$\cdot$m$\cdot$s$^{-2}$), which
    is the natural scale of gravitational coupling.

    \item $\kB = 1.381\times10^{-23}$~J/K is Boltzmann's constant.

    \item $\lnTwo = \ln 2 \approx 0.6931$ is dimensionless --- the conversion
    factor between base-2 (Shannon) and natural logarithm entropy conventions.

    \item $\Sent$ is the \textbf{entanglement entropy density} in
    bit$\cdot$m$^{-3}$ --- the number of bits of quantum entanglement per
    unit volume in the laboratory system. This is not the same as vacuum
    entanglement entropy (which obeys an area law at causal horizons). The
    distinction is critical and is addressed fully in
    Section~\ref{sec:extrapolation}.

    \item $g_{\mu\nu}$ is the \textbf{metric tensor} --- dimensionless, it
    encodes the geometry of spacetime.
\end{itemize}

\noindent The factor $\kappatilde\,c^4/(k_B\lnTwo)$ serves as the effective
coupling between the information-theoretic quantity $\Sent$ (dimensionless
per unit volume) and the geometric quantity $G_{\mu\nu}$ (inverse length
squared). Its dimensional structure is verified explicitly in
Section~\ref{sec:dimensional_rigor}.

The physical entropy density corresponding to $\Sent$ is:
\begin{equation}
    \mathcal{S}_{\text{ent}} = \Sent \cdot \kB\lnTwo
    \quad \left[\frac{\text{J}}{\text{K}\cdot\text{m}^3}\right]
    \label{eq:entropy_density_conversion}
\end{equation}

\noindent This conversion --- multiplying a dimensionless bit count per
volume by Boltzmann's constant and $\ln 2$ --- is what bridges
information theory (bits) to thermodynamics (joules per kelvin).
It is the foundational dimensional step without which the modified
Einstein equation is ill-defined.

% --------------------------------------------------------
\subsection{The Gravitational Source Term for a Perfect Fluid}
\label{subsec:source_term}
% --------------------------------------------------------

In general relativity, the gravitational effect of matter is not determined
by energy density alone but by the combination
$\rho_{\text{grav}} + 3p_{\text{grav}}/c^2$, which appears in the
Raychaudhuri equation governing how a bundle of nearby geodesics (freely
falling paths) converges or diverges. For a perfect fluid with energy density
$\rho$ and pressure $p$:

\begin{itemize}
    \item \textbf{Attractive gravity}: $\rho + 3p/c^2 > 0$
    (ordinary matter: $\rho > 0$, $p \geq 0$, so this is always satisfied)
    \item \textbf{Repulsive gravity}: $\rho + 3p/c^2 < 0$
    (requires sufficiently negative pressure, like dark energy or a
    negative cosmological constant)
\end{itemize}

Incorporating the entanglement entropy source term from
Eq.~\eqref{eq:modified_einstein}, the full gravitational source becomes:

\begin{equation}
    \rho_{\text{grav}} + \frac{3p_{\text{grav}}}{c^2}
    = \underbrace{\rho + \frac{3p}{c^2}}_{\text{ordinary matter}}
    + \underbrace{\frac{3\kappatilde\,c^2}{8\pi G\,\kB\lnTwo}\,
    \Sent}_{\text{entanglement contribution}}
    \label{eq:source_term}
\end{equation}

\noindent When $\kappatilde < 0$ and $\Sent > 0$, the entanglement
contribution is \textbf{negative} --- it reduces the total gravitational
source. For sufficiently high entanglement entropy density with sufficiently
negative $\kappatilde$, the total source can become negative, producing
\textbf{repulsive spacetime curvature}. This is the mechanism by which
high-entanglement quantum systems might generate repulsive gravity without
requiring exotic matter, negative mass, or any physics beyond the Standard
Model and general relativity.

To make this concrete: for the total source to become negative requires:

\begin{equation}
    \Sent > \frac{8\pi G\,\kB\lnTwo}{3|\kappatilde|\,c^2}
    \left(\rho + \frac{3p}{c^2}\right)
    \label{eq:repulsion_threshold}
\end{equation}

This sets the minimum entanglement entropy density required for repulsive
curvature to dominate over ordinary matter attraction. At laboratory scales
with $\kappatilde \sim -10^{-5}$, this threshold is extraordinarily high ---
which is consistent with the absence of observed gravity anomalies in any
current quantum experiment.

% --------------------------------------------------------
\subsection{First-Principles Estimate of $\kappatilde$}
\label{subsec:kappa_derivation}
% --------------------------------------------------------

\begin{framed}
\noindent\textbf{Scope boundary:} The following calculation yields an ideal
value $\kappatilde = -1/4$ under a specific set of assumptions that are
valid at causal horizons (Rindler, black hole event horizons). The extension
of this result to laboratory-scale entanglement volumes requires an additional
physical hypothesis --- the holographic Planck-scale regulator assumption ---
which is stated explicitly at the point where it enters the derivation. The
result should therefore be read as:

\begin{center}
$\kappatilde = -1/4$ \emph{under the holographic Planck-scale regulator
hypothesis}
\end{center}

\noindent not as a fully derived consequence of established physics.
The scientific value of this result lies in providing a specific,
falsifiable prediction --- not in claiming more theoretical authority
than the derivation warrants.
\end{framed}

\subsubsection{Step 1: Jacobson's Thermodynamic Gravity Framework}

Jacobson~\cite{jacobson1995} derived Einstein's equations by applying
thermodynamic reasoning to local causal horizons. The key steps are:

Consider a small patch of spacetime near any point $p$. An observer
accelerating with proper acceleration $a$ through this region perceives
a local Rindler horizon --- a causal boundary that prevents them from
receiving signals from behind. The observer's world is divided into an
accessible region and a causally inaccessible region by this horizon.

The accelerated observer experiences the quantum vacuum as a thermal
state with the \textbf{Unruh temperature}~\cite{unruh1976}:

\begin{equation}
    T_U = \frac{\hbar\,a}{2\pi\,c\,\kB}
    \label{eq:unruh_temperature}
\end{equation}

\noindent \textbf{How to read this:} The numerator $\hbar a$ has units of
energy (joules) --- it is the quantum of action ($\hbar$, in J$\cdot$s)
times the acceleration ($a$, in m$\cdot$s$^{-2}$), giving J$\cdot$m$^{-1}$,
which combined with $c$ in the denominator gives J, and divided by $k_B$
gives kelvin. The factor $2\pi$ is geometric, arising from the Rindler
horizon geometry. Importantly, an inertial observer in the same region
sees the vacuum --- zero particles, zero temperature. The Unruh temperature
is \emph{observer-dependent}, a consequence of the fact that the definition
of ``particle'' depends on the trajectory through spacetime.

For a horizon area element $dA$, the associated Bekenstein-Hawking entropy
is~\cite{bekenstein1973,hawking1975}:

\begin{equation}
    dS_{\mathrm{BH}} = \frac{\kB\,c^3}{4G\hbar}\,dA
    \label{eq:bekenstein_hawking}
\end{equation}

\noindent \textbf{How to read this:} The combination $c^3/(G\hbar)$ has
units of m$^{-2}$, so $dS_{\mathrm{BH}}$ has units of J/K per m$^2$,
i.e., entropy per unit area. The famous ``$1/4$'' prefactor is a
result from quantum field theory in curved spacetime that took decades
to establish. It says that each Planck area $\ell_P^2 = \hbar G/c^3$
of horizon carries $\kB/4$ of entropy --- a universal result independent
of the details of the matter content.

The Clausius relation $\delta Q = T\,dS$, applied to energy $\delta Q$
flowing across this horizon area element, gives:

\begin{equation}
    \delta Q_{\mathrm{BH}} = T_U\,dS_{\mathrm{BH}}
    = \frac{\hbar\,a}{2\pi\,c\,\kB} \cdot \frac{\kB\,c^3}{4G\hbar}\,dA
    = \frac{a\,c^2}{8\pi G}\,dA
    \label{eq:clausius_bh}
\end{equation}

\noindent By identifying this heat flux with the energy-momentum flux
$T_{\mu\nu}k^\mu d\Sigma^\nu$ across the horizon (where $k^\mu$ is the
horizon-generating Killing vector and $d\Sigma^\nu$ is the area element),
Jacobson recovers the Einstein equations. This is the thermodynamic
derivation of GR.

\subsubsection{Step 2: Entanglement Entropy Contribution to the Clausius
Relation}

If the quantum state of matter has entanglement entropy density $\Sent$
(bit$\cdot$m$^{-3}$), the corresponding physical entropy density is
$\mathcal{S}_{\mathrm{ent}} = \Sent\cdot\kB\lnTwo$.

The entanglement entropy contribution to the Clausius relation across a
horizon area element $dA$ is:

\begin{equation}
    dS_{\mathrm{ent}} = \frac{\mathcal{S}_{\mathrm{ent}}}{\kB}\cdot
    \frac{dV}{4\ell_P}
    \label{eq:dS_ent}
\end{equation}

\noindent where $dV = \ell_P\,dA$ is the volume element immediately behind
the horizon (one Planck length deep) and
$\ell_P = \sqrt{\hbar G/c^3} \approx 1.616\times10^{-35}$~m is the
Planck length.

\begin{framed}
\noindent\textbf{[Holographic Planck-scale regulator hypothesis]:}
The assignment $dV = \ell_P\,dA$ is the critical assumption in this
derivation. It treats the Planck length as the natural depth scale behind
a horizon over which entanglement entropy accumulates. This is motivated by:
\begin{enumerate}
    \item The Bekenstein-Hawking result, which assigns one bit of entropy
    per $4\ell_P^2$ of horizon area --- implying a one-Planck-length
    correlation depth
    \item Holographic principles~\cite{susskind1995} suggesting that all
    information in a volume can be encoded on its boundary at Planck
    density
    \item The ultraviolet structure of quantum field theory, which is
    regulated at the Planck scale
\end{enumerate}
This assumption is valid at causal horizons. Its extension to laboratory
volumes --- where no causal horizon exists --- is an extrapolation
hypothesis, not a derivation. See Section~\ref{sec:extrapolation} for the
complete statement of what is assumed versus proven.
\end{framed}

Substituting $\mathcal{S}_{\mathrm{ent}} = \Sent\cdot\kB\lnTwo$ and
$dV = \ell_P\,dA$ into Eq.~\eqref{eq:dS_ent}:

\begin{equation}
    dS_{\mathrm{ent}} = \frac{\Sent\cdot\kB\lnTwo}{\kB}\cdot
    \frac{\ell_P\,dA}{4\ell_P}
    = \frac{\Sent\cdot\lnTwo}{4}\,dA
    \label{eq:dS_ent_simplified}
\end{equation}

\noindent Note that the Planck length $\ell_P$ cancels exactly in the
ratio $dV/\ell_P = dA$ --- this is not a coincidence but a consequence
of the holographic regulator assumption, which is precisely designed to
make $\ell_P$ cancel and leave a finite result.

\subsubsection{Step 3: Modified Clausius Relation}

The total heat flux across the horizon area element, including the
entanglement entropy contribution, is:

\begin{align}
    \delta Q_{\mathrm{eff}} &= T_U\,\left(dS_{\mathrm{BH}}
    + dS_{\mathrm{ent}}\right) \nonumber \\
    &= T_U\,dS_{\mathrm{BH}}
    + \frac{\hbar\,a}{2\pi\,c\,\kB}\cdot
    \frac{\Sent\cdot\kB\lnTwo}{4}\,dA \nonumber \\
    &= \frac{a\,c^2}{8\pi G}\,dA
    + \frac{\hbar\,a\,\Sent\,\lnTwo}{8\pi\,c}\,dA
    \label{eq:modified_clausius}
\end{align}

\noindent The first term on the right recovers Jacobson's original result
(Eq.~\eqref{eq:clausius_bh}). The second term is the new entanglement
entropy contribution.

\subsubsection{Step 4: Identification of Effective Stress-Energy}

The additional heat flux in Eq.~\eqref{eq:modified_clausius} acts as an
effective energy-momentum contribution across the horizon. Identifying
$\delta Q_{\mathrm{eff,new}} = T^{\mathrm{eff}}_{\mu\nu}\,k^\mu\,
d\Sigma^\nu$ and using the surface gravity relation $a = c^2\kappa_s$
(where $\kappa_s$ is the surface gravity at the horizon, not to be confused
with the coupling constant $\kappa$), and following the same geometric steps
as Jacobson's original derivation:

\begin{equation}
    T^{\mathrm{eff}}_{\mu\nu} = -\frac{c^4}{32\pi G}\,\Sent\,g_{\mu\nu}
    \label{eq:T_eff}
\end{equation}

\noindent \textbf{How to read this:} The negative sign arises because the
entanglement entropy contribution acts as an effective negative pressure ---
it tends to push geodesics apart rather than together. The combination
$c^4/(32\pi G)$ has units of Pa$\cdot$m$^3$/bit (pressure times volume per
bit), which when multiplied by $\Sent$ (bit/m$^3$) gives Pa = kg/(m$\cdot$s$^2$),
consistent with the stress-energy tensor. The factor of $g_{\mu\nu}$
means this is an isotropic contribution --- equal in all directions, like
a cosmological constant term.

\subsubsection{Step 5: Reading Off the Coupling Constant}

Comparing Eq.~\eqref{eq:T_eff} with the entanglement source term in
Eq.~\eqref{eq:modified_einstein}, we identify the relationship:

\begin{equation}
    \kappatilde\,\frac{c^4}{\kB\lnTwo}\,\Sent
    = 8\pi G \cdot T^{\mathrm{eff}}_{\mu\nu}/g_{\mu\nu}
    = 8\pi G \cdot \left(-\frac{c^4}{32\pi G}\,\Sent\right)
    = -\frac{c^4}{4}\,\Sent
    \label{eq:kappa_comparison}
\end{equation}

Solving for $\kappatilde$:

\begin{equation}
    \kappatilde\,\frac{c^4}{\kB\lnTwo} = -\frac{c^4}{4}
    \quad \Longrightarrow \quad
    \kappatilde = -\frac{\kB\lnTwo}{4}
    \label{eq:kappa_intermediate}
\end{equation}

Wait --- this gives $\kappatilde = -k_B\ln2/4$, which is not dimensionless.
We must resolve this. The issue is that $\Sent$ in the effective stress-energy
tensor Eq.~\eqref{eq:T_eff} is the \emph{physical} entropy density
(J/(K$\cdot$m$^3$)), while $\Sent$ in the modified Einstein equation
Eq.~\eqref{eq:modified_einstein} is the \emph{information} entropy density
(bit/m$^3$). Converting using
$\mathcal{S}_{\mathrm{ent}} = \Sent\cdot\kB\lnTwo$:

\begin{equation}
    T^{\mathrm{eff}}_{\mu\nu} = -\frac{c^4}{32\pi G}
    \cdot\frac{\mathcal{S}_{\mathrm{ent}}}{\kB\lnTwo}\cdot g_{\mu\nu}
    = -\frac{c^4}{32\pi G}\,\Sent\,g_{\mu\nu}
\end{equation}

\noindent where in the last step we used $\Sent = \mathcal{S}_{\mathrm{ent}}
/(\kB\lnTwo)$ (converting physical entropy density back to bit density).
Now comparing with the entanglement source term in
Eq.~\eqref{eq:modified_einstein}:

\begin{equation}
    \frac{\kappatilde\,c^4}{\kB\lnTwo}\cdot\Sent
    \stackrel{!}{=} 8\pi G\cdot\left(-\frac{c^4}{32\pi G}\,\Sent\right)
    = -\frac{c^4\,\Sent}{4}
\end{equation}

\noindent Dividing both sides by $c^4\,\Sent/(\kB\lnTwo)$:

\begin{equation}
    \boxed{\kappatilde = -\frac{1}{4}}
    \label{eq:kappa_result}
\end{equation}

\noindent This is the \textbf{ideal coupling constant} in the absence
of environmental decoherence, under the holographic Planck-scale
regulator hypothesis. It is dimensionless, as required.

The negative sign confirms that high entanglement entropy density produces
effective negative pressure, consistent with the repulsive curvature
interpretation. The magnitude $1/4$ is a geometric factor inherited from
the Bekenstein-Hawking ``1/4'' in the area-entropy relation.

% --------------------------------------------------------
\subsection{Environmental Screening}
\label{subsec:screening}
% --------------------------------------------------------

The ideal value $\kappatilde = -1/4$ applies to a perfectly isolated,
maximally coherent quantum system at zero temperature --- a physical
idealization. In any realistic laboratory system, environmental
decoherence suppresses the long-range quantum correlations that drive the
entanglement-gravity coupling. The realistic coupling is:

\begin{equation}
    \kappatilde = -\frac{1}{4}\,\alphascreen
    \label{eq:realistic_kappa}
\end{equation}

\noindent where $\alphascreen \in [0,1]$ is the \textbf{environmental
screening factor} --- the fraction of the ideal entanglement-gravity
coupling that survives decoherence.

A phenomenological model for $\alphascreen$ combines two physical
suppression mechanisms:

\begin{equation}
    \alphascreen = \frac{1}{1 + (\Gamma\tau_D)^2}
    \cdot e^{-(R/\xi)^2}
    \label{eq:alpha_screen}
\end{equation}

\noindent where each factor has the following physical meaning:

\begin{itemize}
    \item \textbf{Temporal suppression} $[1 + (\Gamma\tau_D)^2]^{-1}$:
    When the product of decoherence rate $\Gamma$ (s$^{-1}$) and the
    characteristic coherence time $\tau_D$ (s) satisfies
    $\Gamma\tau_D \gg 1$, the system decoheres faster than it can maintain
    coherence, and $\alphascreen \to 0$. When $\Gamma\tau_D \ll 1$, the
    system is near-perfectly coherent and $\alphascreen \to 1$. The
    Lorentzian form is motivated by Kubo-Martin-Schwinger spectral theory
    for weakly coupled open systems.

    \item \textbf{Spatial suppression} $e^{-(R/\xi)^2}$: Long-range
    quantum correlations decay on a characteristic coherence length $\xi$
    (m). When the system size $R \gg \xi$, spatial decoherence destroys
    the long-range entanglement and $\alphascreen \to 0$. The Gaussian
    form is standard for many-body decoherence in the thermodynamic limit.

    \item $\Gamma$: Decoherence rate, s$^{-1}$. For trapped $^{87}$Rb
    atoms at $1\,\mu$K: $\Gamma \sim 1$--$10$~s$^{-1}$.

    \item $\tau_D$: Characteristic decoherence time, s. Distinct from
    $1/\Gamma$: $\tau_D$ is the coherence time of the macroscopic
    quantum state; $\Gamma$ is the rate of individual decoherence events.

    \item $R$: Linear size of the entangled system, m.

    \item $\xi$: Coherence length of the entangled state, m. Depends on
    temperature, trapping geometry, and entanglement protocol.
\end{itemize}

\noindent\textbf{Important caveat:} Eq.~\eqref{eq:alpha_screen} is a
\emph{phenomenological ansatz}. The specific functional forms (Lorentzian
in $\Gamma\tau_D$, Gaussian in $R/\xi$) are physically motivated but not
derived from a Lindblad master equation for any specific system. A
rigorous derivation of $\alphascreen$ for trapped $^{87}$Rb GHZ states
from the open-quantum-system Lindblad equation is an \textbf{[OPEN PROBLEM]}
in this research program.

Numerical estimates for trapped $^{87}$Rb atoms at $1\,\mu$K yield:

\begin{equation}
    \alphascreen \in [10^{-4},\;10^{-2}]
    \label{eq:alpha_range}
\end{equation}

\noindent giving a realistic coupling range:

\begin{equation}
    \kappatilde \in [-2.5\times10^{-3},\;-2.5\times10^{-5}]
    \label{eq:kappa_range}
\end{equation}

This range is entirely consistent with existing experimental null results,
which constrain $|\kappatilde| < 10^{-10}$ (see Table~\ref{tab:bounds}).
The gap between the current best experimental bound ($10^{-10}$) and the
predicted realistic coupling ($\sim 10^{-4}$ to $10^{-5}$) means this
framework predicts a \emph{detectable signal} at proposed sensitivity
levels --- not a signal that is guaranteed to be beyond reach.

% --------------------------------------------------------
\subsection{Conservation of Energy-Momentum}
\label{subsec:conservation}
% --------------------------------------------------------

In standard general relativity, the Bianchi identity $\nabla^\mu G_{\mu\nu}
= 0$ implies that the stress-energy tensor is independently conserved:
$\nabla^\mu T_{\mu\nu} = 0$. This is the covariant expression of local
energy-momentum conservation, tested to extraordinary precision in the
solar system and in precision atomic physics.

With the entanglement entropy source term, the modified Einstein equation
Eq.~\eqref{eq:modified_einstein} implies instead:

\begin{equation}
    \nabla^\mu T_{\mu\nu} = -\frac{1}{8\pi G}\,\partial_\nu
    \left(\kappatilde\,\frac{c^4}{\kB\lnTwo}\,\Sent\right)
    \label{eq:conservation_modified}
\end{equation}

\noindent This means that ordinary matter stress-energy is \emph{not}
independently conserved when the entanglement entropy density has a
spatial gradient. The right-hand side represents an energy-momentum
exchange between ordinary matter and the entanglement sector --- mediated
by the spatial gradient of $\Sent$.

\textbf{Tension with precision tests:} Local energy-momentum conservation
has been verified to parts-per-billion precision in Solar System dynamics
and to parts-per-million in precision atomic clocks. This places a very
stringent constraint: the right-hand side of Eq.~\eqref{eq:conservation_modified}
must be negligibly small in all tested environments.

The natural estimate of this term is:

\begin{equation}
    \left|\frac{1}{8\pi G}\,\partial_\nu\left(\kappatilde\,
    \frac{c^4}{\kB\lnTwo}\,\Sent\right)\right|
    \sim \frac{|\kappatilde|\,c^4\,\Sent}{8\pi G\,\kB\lnTwo\,L}
\end{equation}

\noindent where $L$ is the spatial scale over which $\Sent$ varies.
With $|\kappatilde| < 10^{-10}$ (current experimental bound) and
$\Sent \sim 10^{20}$~bit/m$^3$ (typical laboratory quantum system),
this is of order $10^{-30}$~Pa/m --- far below any tested energy-momentum
flux. Conservation is therefore consistent with current bounds.

\textbf{Prediction:} Any experimental detection of nonzero $\kappatilde$
implies a corresponding violation of standard energy-momentum conservation,
at the level $\sim |\kappatilde|\cdot c^4 \cdot\Sent / (8\pi G\kB\lnTwo L)$.
This is itself a falsifiable prediction of the framework.
